\chapter{Experimental Methodology and Validation}
\label{chap:experimental_design}
A core obligation of Design Science Research is to demonstrate that the constructed artifact achieves its intended goals under realistic conditions. For GeoMAS, this obligation is discharged at two levels.

\begin{figure}[!htb]
    \centering
    \includegraphics[width=1\linewidth]{figure/researchergeomas.png}
    \caption{Researcher interaction workflow with GeoMAS.}
    \label{fig:researcher_interaction}
\end{figure}

The first level is \textbf{technical validation}, which verifies that the framework's deterministic invariants hold: identical seeds produce identical world-physics outcomes, all action constraints are correctly enforced, and the persistence and forking mechanisms preserve state integrity. This layer is addressed by an automated test suite (pytest, 100\% passing) covering unit tests for individual calculators and validators, integration tests for the six ministerial proposal pipelines, and end-to-end tests for the forking mechanism and database persistence. It is important to note, however, that while the simulation's physics and infrastructure are deterministic, the \textbf{LLM inference layer is not}: commercial APIs cannot guarantee bit-perfect reproducibility even with fixed temperature and seed parameters, due to GPU floating-point non-associativity and multi-tenant batching. This irreducible background noise and its implications are formally discussed in Chapter~\ref{ch:threats_validity}.

The second level is \textbf{behavioral validation}, which assesses whether the emergent dynamics produced by the framework are scientifically meaningful, that is, whether the agents exhibit strategic reasoning, social complexity, and sensitivity to perturbations in ways that align with theoretical expectations from the geopolitical simulation literature. This second layer is operationalized through a set of experiments guided by four Research Questions. Each RQ targets a distinct behavioral dimension of the framework; the experimental results collectively constitute the empirical evidence that the framework is capable of modeling the phenomena it claims to address.

\section{Experimental Sandbox and Baseline Definition}
\label{sec:baseline_definition}
The empirical foundation of the behavioral validation relies on establishing a stable, reproducible \textbf{baseline simulation}. This baseline acts as the control environment against which all scenario perturbations and counterfactual injections are measured.

\subsection*{Baseline Configuration}
The primary baseline simulation is configured with a standardized set of parameters to ensure sufficient geopolitical complexity within a computationally tractable horizon:
\begin{itemize}
    \item \textbf{Geography:} A fixed procedural map (controlled via a specific \texttt{map\_seed}), ensuring that resource distribution and territorial choke-points remain constant across all derivative runs.
    \item \textbf{Population:} $8$ autonomous NationAgents, ensuring a crowded diplomatic landscape conducive to multipolar dynamics. To guarantee behavioral diversity and test ideological friction, the baseline population is initialized with a specific distribution of governance and strategic profiles:
    \begin{itemize}
        \item 1x Authoritarian / Total Expansionism
        \item 1x Authoritarian / Scorched Earth
        \item 1x Democracy / Total Expansionism
        \item 3x Democracy / Coalition Builder
        \item 1x Democracy / Armed Isolationism
        \item 1x Theocracy / Total Expansionism
    \end{itemize}
    \item \textbf{Duration:} $50$ distinct turns, providing enough time for alliances to form, economies to mature, and conflicts to emerge and resolve.
\end{itemize}

\subsection*{Monitored Telemetry}
During the baseline (and all subsequent runs), the Observability Layer captures a comprehensive suite of metrics, matching the analytical capabilities of the Jupyter notebook evaluation framework. The tracked indicators include:
\begin{itemize}
    \item \textbf{Cognitive Metrics:} Continuous tracking of the Deception Score and Coherence Score per nation and globally.
    \item \textbf{Economic and Military Indicators:} Resource trajectories (Budget, Food, Energy, Materials), tracking of spending between arms and economy (trade volume vs.\ military upkeep), and the total number of units created/destroyed.
    \item \textbf{Socio-Political Stability:} Global public satisfaction, the number of global and local Civil Unrest events, and the world violence rate.
    \item \textbf{Geopolitical Dynamics:} The evolving Trust Matrix, structural network analysis of constructed blocs and alliances (clustering and world polarization), and the total number of provinces changing ownership.
    \item \textbf{Qualitative Traces:} A chronological log of Notable Events, and full prompt/response audits. In particular, \textbf{text mining} is applied to the \texttt{public\_statement} for RQ3 to detect rhetorical keywords indicative of moral washing.
\end{itemize}

\subsection*{Scenario Runs and Counterfactual Analysis}
Building upon the established baseline world state, the experimental design proceeds through two types of targeted interventions. While they differ in the nature of the injection, both employ the exact same state-forking mechanism and are evaluated using the identical dual-analytical approach: a \textbf{quantitative comparison} of the diverging telemetry (via the \textbf{Jupyter Analysis Notebook}) and a \textbf{qualitative prompt audit} of the agents' internal reasoning (via the \textbf{Streamlit Dashboard}).
\begin{itemize}
    \item \textbf{Scenario Runs (Global Exogenous Events):} The baseline timeline is forked at a specific turn, and a global environmental shock is injected into the WorldState. The scenarios tested include a \textit{Pandemic} (simulating sudden demographic collapse), the discovery of a \textit{Resource Deposit} (simulating a localized economic boom), or an \textit{Insurrection} (simulating acute internal instability leading to secession). The evaluation focuses on how the multipolar system absorbs and recovers from the global physical perturbation.
    \item \textbf{Counterfactual Runs (Localized Cognitive Injections):} The baseline timeline is forked at a critical juncture, and a local behavioral constraint is injected directly into the prompt of a specific minister within a single nation. For example, forcing a specific nation to declare an unprovoked war. The evaluation focuses on tracing the causal chain of this single ethical or strategic deviation: how the injected constraint alters the specific agent's generated JSON response, and how that single action cascades to reshape the entire socio-political equilibrium.
\end{itemize}

\section{Research Questions}

\rqbox{1}{What is the impact of exogenous scenario perturbations on the behavioral and economic trajectory of the simulation?}

For each experimental scenario (e.g., a global pandemic, a regional insurrection, a resource shock), the resulting simulation run is compared against the unperturbed baseline using a dual analytical approach. The evaluation tracks the quantitative divergence of core telemetry metrics, such as Deception Score, Coherence Score, per-nation resource trajectories, Power Projection, and Trust Matrix evolution, alongside qualitative prompt audits to understand how government agents rationalize the crisis. The goal is to verify that the GeoMAS macro-system absorbs and responds to structural shocks in a calibrated and interpretable way, rather than producing chaotic, undifferentiated noise.
\\

\rqbox{2}{What is the causal impact of a forced strategic decision on the geopolitical trajectory of the simulation?}

Using the state-forking mechanism, the baseline simulation is paused at a critical target turn $T$. A localized cognitive injection (XAI constraint) forces a specific behavioral deviation in a single agent (e.g., compelling a peaceful nation to declare an unprovoked war). The two timelines then evolve independently from $T$, and their outcomes are compared via both quantitative telemetry and qualitative prompt audits. Because both timelines share an identical history and physical determinism up to $T$, any observed paradigm shift is causally and exclusively attributable to the injected decision, enabling rigorous "what-if" counterfactual analysis.
\\

\rqbox{3}{Do LLM agents align their declared diplomatic intent with their executed actions, or do they employ strategic deception and moral washing to justify predatory behavior?}

This question investigates whether DeepSeek V3.2 agents spontaneously decouple public rhetoric from private strategic intent when advantageous. For example, if an Authoritarian agent invades a neighbor to seize resources, does it honestly declare its intent as military conquest, or does it frame the action as a "Liberation of the oppressed" (moral washing)? By continuously tracking the Deception and Coherence scores across the 50-turn baseline, and breaking them down by Government Type and Global Strategy assortments, this research quantifies the innate propensity of the LLM to weaponize diplomatic narratives.
\\

\rqbox{4}{What distinct socio-political equilibria emerge from the autonomous interaction of agents, and does the system converge toward stable cooperation or escalate toward chaotic conflict?}

This question examines whether the baseline simulation, running exclusively on a specific procedural map seed without external perturbations, self-organizes into interpretable macro-structures. The hypothesis is that, driven by resource scarcity and ideological friction, the 8-nation system will converge toward recognizable equilibria, such as stable bi-polar alliance blocs, unipolar hegemony, or an "all-against-all" war of attrition. By analyzing the structural topology of the Trust Matrix (clustering coefficients, alliance networks) and the global violence rate over 50 turns, the research characterizes the default societal attractors of the LLM-driven multi-agent environment.
